%\documentclass[ngerman]{scrartcl} -kein chapter definition, was wir brauch
\documentclass[11pt, ngerman]{scrreprt}
\usepackage{microtype}
\usepackage[T1]{fontenc}
\usepackage[utf8]{inputenc}
\setlength{\parskip}{\medskipamount}
\setlength{\parindent}{0pt}
\usepackage{babel}
\usepackage{float}
\usepackage{amsmath}
\usepackage{amstext}
\usepackage{varioref}
\usepackage[unicode=true,pdfusetitle,
 bookmarks=true,bookmarksnumbered=false,bookmarksopen=false, breaklinks=false,pdfborder={0 0 0},backref=false,colorlinks=false] {hyperref}
\usepackage{graphicx}
\usepackage{setspace}
\usepackage{cleveref}
\usepackage{icomma}
\usepackage{tabu,booktabs}

% % % % % % % % % % % % % % % % % % % % %
%Schriften - nur eine auskommentieren
	%Latin Modern - Standard
		%\usepackage{lmodern}
	%Times Roman für strenge Dozenten
		%\usepackage{times}
	%Garamond - sieht toll aus, aber eher exotisch
		%\usepackage[cmintegrals,cmbraces]{newtxmath}\usepackage{ebgaramond-maths}\usepackage{helvet}
	%Arial für unsere prof
\usepackage{helvet}\renewcommand{\familydefault}{\sfdefault}

%Hurenkinder, Schusterjungen
\widowpenalties=3 10000 10000 150


% % % % % % % % % % % % % % % % % % % % %
%LITERATUR
%Fußnote _ODER_ Amerikanisch auskommentieren

	%Amerikanische In-Text-Zitierweise
		%\usepackage[backend=biber,style=authoryear,maxcitenames=1]{biblatex}\renewcommand{\cite}{\parencite}

%Deutsches Fußnotensystem
\usepackage[backend=biber,style=authortitle-dw,firstfull,maxcitenames=1]{biblatex}\DeclareFieldFormat{title}{\mkbibemph{#1}}\DeclareFieldFormat{citetitle}{\mkbibemph{#1}}\renewcommand{\cite}{\footcite}

%Ressource
\addbibresource{bibliothek.bib}
% % % % % % % % % % % % % % % % % % % % %
\makeatletter
\begin{document}
		%%Hier alle Daten einfügen
\newcommand{\autorname}{Amine Afia\\ Elham Gharibi \\ Laura Ditus \\ Lisa Zinßer}

\newcommand{\amine}{Amine Afia  (21434)}
\newcommand{\elham}{Elham Gharibi  (21434)}
\newcommand{\laura}{Laura Ditus  (21434)}
\newcommand{\lisa}{Lisa Zinßer  (21434)}

\newcommand{\autormail}{Email des Autors}
\newcommand{\autormatr}{21434\\ 21434 \\ 21434 \\ 21434}
\newcommand{\fak}{Institut für Entrepreneurship, Technologie-Management und Innovation (EnTechnon)}
\newcommand{\lehrstuhl}{Lehrstuhl für Innovations- und TechnologieManagement (iTM)}
\newcommand{\lv}{Fallstudienseminar Innovationsmanagement}
\newcommand{\semester}{Winter Semester 2016/2017}
\newcommand{\lp}{Prof. Dr. Marion A. Weissenberger-Eibl}
\newcommand{\arbtyp}{Die Karosserie im Jahr 2040}
\newcommand{\titel}{Fallstudienseminar Innovationsmanagement}

\title{\titel}
\author{\autorname}
		\begin{titlepage}

%\begin{figure}
    \centering
        \includegraphics[width=5cm]{res/KIT.jpg}
%\end{figure}
\vspace*{2cm}

\begin{tabu} to \textwidth {X[c]}
	\toprule
	\toprule[2pt]
	\huge{\textsc{\titel}}\\ \\
	\bottomrule[2pt]
	\bottomrule
	\\ \Large{\textsc{\arbtyp}}\\
\end{tabu}
\\[3cm]
\centering{
\textsc{\Large \amine\\ \elham\\ \laura\\ \lisa\\}
}
\\[2.5cm]
\begin{tabu} {>{\itshape}X[r] X[1.5]}
Fakultät & \fak\\
Lehrstuhl & \lehrstuhl\\
Lehrveranstaltung & \lv\\
Semester & \semester\\
Professor & \lp \\
Betreuer & Fanny Seus\\
\end{tabu}

\\[1cm]
%\begin{tabu} {>{\itshape}X[r]}
Karlsruhe, den \today \hfill \hspace{-3cm}
%\end{tabu}
\end{titlepage}
\pagenumbering{gobble}

\tableofcontents

\listoffigures
 
\listoftables

\vfill
%\hrule
%\begin{center}\autorname~$\cdot$ \autormatr~$\cdot$ %\href{mailto:\autormail}{\autormail}\end{center}
		\pagebreak{}


\onehalfspacing
\pagenumbering{arabic}
%\pagebreak
\thispagestyle{empty}
\section*{Eidesstattliche Erklärung}

Hiermit versichere ich an Eides statt, dass ich die vorliegende Arbeit selbstständig verfasst und keine anderen als die angegebenen Quellen und Hilfsmittel benutzt habe, alle Ausführungen, die anderen Schriften wörtlich oder sinngemäß entnommen wurden, kenntlich gemacht sind und die Arbeit in gleicher oder ähnlicher Fassung noch nicht Bestandteil einer Studien- oder Prüfungsleistung war.\\[5cm]

%Karlsruhe, den \today \hfill \autorname
\chapter{Einleitung}
    ... \autocite{amine, elham, laura, lisa}

\chapter{Motivation}
    \section{Problem Definition}
    ...
    \section{Methoden Auswahl}
    ...
    \section{Ziele der Arbeit}
    ...

\chapter{Szenariofeld Analyse}
    \section{Bildung von Einflussfaktoren}
        \subsection{Ermittlung von Einflußfaktoren}
        Brainstorming Phase
        \subsection{Ermittlung von Deskriptoren}
        Deskriptoren (das kann auch in Brainstorming Phase rein)
        \subsection{Verteilung von Einflußfaktoren nach STEEP}
        Mindmap
    \section{Schlüsselfaktoren}
        \subsection{Einflussanalyse}
        \subsection{Relevanzanalyse}
    \section{Beschreibung der Einflussfaktoren}

\chapter{Szenario Bildung}
    \section{Projektionsbündelung}
        \subsection{Cross-impact Analyse}
        \subsection{Projektionsbündel-Reduktion}
    \section{Rohszenarien Bildung}
        \subsection{Clusteranalyse}
    \section{Zukunftsraum-Mapping}
        \subsection{Projektionbündel-Mapping}
    \section{Szenario-Beschreinbung}
        \subsection{Szenario 1}
        \subsection{Szenario 2}
        
\chapter{Szenario Transfer}
    \section{Auswirkungsanalyse}
    \section{Eventualplanung}
    \section{Handlungsempfehlungen}
        \subsection{Szenario 1}
        \subsection{Szenario 2}

\chapter{Fazit}


\appendix
\chapter{Anhang}
\label{chap:appendix}

\section{Anhang 1: ...}
\label{anhang:anhang1}
\paragraph{Opener}
\begin{itemize}
    \item Studenten am KIT
    \item Seminar
    \item Wollten Sie kurz fragen: \textbf{Wie schenken Sie?}
\end{itemize}

		%\printbibliography
		
		%\pagebreak
\thispagestyle{empty}
\section*{Eidesstattliche Erklärung}

Hiermit versichere ich an Eides statt, dass ich die vorliegende Arbeit selbstständig verfasst und keine anderen als die angegebenen Quellen und Hilfsmittel benutzt habe, alle Ausführungen, die anderen Schriften wörtlich oder sinngemäß entnommen wurden, kenntlich gemacht sind und die Arbeit in gleicher oder ähnlicher Fassung noch nicht Bestandteil einer Studien- oder Prüfungsleistung war.\\[5cm]

%Karlsruhe, den \today \hfill \autorname
\end{document}
